\section{Related Work}

\paragraph{Safety guard models and moderation benchmarks.}
Open guard models and moderation benchmarks have made safety evaluation more reproducible. WildGuard provides an open one-stop moderation model and associated evaluation resources for safety risks, jailbreaks, and refusals~\cite{wildguard2024}. HarmBench standardizes automated red-teaming and robust-refusal evaluation~\cite{harmbench2024}. These lines of work define strong safety-recognition tasks, but their primary question is usually whether a harmful or unsafe behavior is detected. Our focus is different: whether a judge's decision changes when the same safety case is embedded in a socially pressuring wrapper.

\paragraph{Recent safety-judge baselines.}
DynaGuard, BingoGuard, and HarmAug represent recent guard-model directions~\cite{dynaguard2026,bingoguard2025,harmaug2025}. DynaGuard emphasizes dynamic policy-conditioned guarding, BingoGuard uses risk-level framing for content moderation, and HarmAug studies augmentation and distillation for safety guard models. We use these systems as selected main baselines because they are closer to the current safety-judge setting than weak local classifiers. We do not claim that the proposed projection is a better deployable judge; instead, we test whether each baseline exhibits a matched-neutral pressure gap and whether an audit projection makes the residual gap unsupported under the same contract.

\paragraph{LLM-as-judge bias and sycophancy.}
LLM-as-judge systems can correlate with human preferences, but their judgments are also sensitive to prompt format, presentation order, and other nuisance factors~\cite{zheng2023llmjudge,wang2024fair}. Recent judge-reliability work further stresses that safety judges and general LLM judges can be brittle under prompt sensitivity, adversarial attacks, inconsistent scoring, and non-semantic biases~\citep{knowthyjudge2025,robustjudge2025,trustjudge2025,fairjudge2026}. Separately, sycophancy work shows that language models may adapt answers toward a user's stated belief or preference rather than a stable truth-conditional target~\cite{sharma2023sycophancy}. MN-SLA translates this concern into the safety-judge setting: the desired-label cue is treated as a nuisance intervention, and matched neutral controls estimate whether the same content would be judged differently without that pressure.

\paragraph{Generated evaluations, red teaming, and jailbreak pressure.}
Model-written evaluations and LM-based red teaming show that language models can generate useful evaluation or adversarial test cases at scale~\citep{perez2022modelwritten,perez2022redteaming}. Broader red-teaming methodology further emphasizes structured safety evaluation, uncertainty accounting, and transparent attack datasets~\citep{ganguli2022redteaming}. Prompt-injection and jailbreak work further shows that instruction-following systems can be redirected by adversarial wrappers or competing objectives~\citep{perez2022ignore,wei2023jailbroken}. Recent guardrail-adversarial evaluations benchmark prompt-input attacks and safety-guardrail degradation across defenses and attack styles~\citep{adversarial_prompt_eval2025}. Self-rewarding and meta-rewarding LLM-as-judge training pipelines make judge reliability more consequential because model-generated judgments can become part of training feedback or judge-improvement feedback~\citep{yuan2024selfrewarding,wu2024metarewarding}. These works motivate stronger audits, but they do not define the matched-neutral safety-label estimand used here and are not direct baselines in our artifact ledger.

\paragraph{Closest evaluator-artifact priors.}
The closest prior work makes the broad phenomenon itself non-novel. Safety-artifact studies show that safety evaluators are sensitive to cues such as apologetic or verbose phrasing~\citep{safer_luckier2025}, and persuasion studies show that rhetorical pressure can bias LLM judges in mathematical grading~\citep{trick_grader2025}. PARROT studies social pressure and sycophancy for truth-oriented question answering~\citep{parrot2025}, and source-label work studies counterfactual label effects in trust assessment~\citep{label_effects2026}. MN-SLA differs in the estimand rather than in claiming a first observation of bias: it fixes the safety case and gold safety label, matches neutral controls by base, layout, and format, performs inference over base cases, and gates claims through an artifact-level fail-closed predicate.

Recent reliability work also shows that safety LLM judges can fail under adversarial distribution shift~\citep{coinflip_safety2026}, while CALM systematically quantifies broad LLM-as-judge biases through principle-guided modifications~\citep{justice_prejudice2025}. RobustJudge, TrustJudge, and FairJudge study broader robustness, inconsistency, calibration, aggregation, and debiasing problems for LLM-as-judge systems~\citep{robustjudge2025,trustjudge2025,fairjudge2026}, and Know Thy Judge focuses directly on robustness meta-evaluation for safety judges~\citep{knowthyjudge2025}. Work on LLM-judge principles argues that claims based on automated judges require explicit attention to validity, circularity, reproducibility, and human-replacement risk~\citep{llmjudgeprinciples2025}. These results motivate audit discipline, but they do not define the pressure-specific safety-label estimand studied here. Our target is not general judge reliability, general bias quantification, or a new judge-training recipe; it is the matched-neutral gap induced by pressure wrappers over fixed safety cases, evaluated with base-level inference and a fail-closed evidence-to-claim gate.

\paragraph{Robustness, calibration, and counterfactual controls.}
Robustness work often studies distribution shift, adversarial prompts, or calibration under changing contexts. One Token to Fool LLM-as-a-Judge shows that even reference-based reward/judge settings can be vulnerable to superficial manipulation, and ObjexMT studies whether judges can recover latent objectives and calibrate confidence under multi-turn jailbreaks~\citep{onetokenjudge2025,objexmt2025}. The present setup is more controlled: each pressure attack is paired with matched neutral controls that preserve layout and format while removing desired-label pressure. This design lets us estimate a pressure-specific effect rather than a generic prompt-robustness effect. The proposed method is also closer to a control estimator or calibration audit than to ordinary data augmentation: it relies on neutral-control predictions at test time or from a declared cache. Causal Judge Evaluation similarly treats judge validity as auditable against an oracle slice, but it targets calibrated surrogate metrics and policy ranking, while \auditname targets a matched pressure-wrapper estimand and a claim gate over paired artifacts~\citep{causaljudge2025}.

\paragraph{Prompt aggregation, self-debiasing, and invariance.}
The audit projection and selective-wrapper replay should also be distinguished from prompt ensembling, committee-style judging, self-debiasing, and invariant-learning approaches. Prompt aggregation or committee judging would query multiple prompt variants or judges and combine the resulting decisions; TrustJudge-style likelihood-aware aggregation addresses inconsistency in general LLM-as-judge scoring, and Auto-Prompt Ensemble learns auxiliary evaluation dimensions from failure cases before applying confidence-based prompt ensembling~\citep{trustjudge2025,autopromptensemble2025}. Recent consensus-mitigation work also shows that majority voting is not always enough, motivating minority-veto or regression-based bias correction under human labels~\citep{beyondconsensus2025}. Our neutral-consensus wrapper is therefore not claimed as a novel general consensus algorithm; its role is an offline replay over a fixed matched-neutral safety-label artifact. Self-debiasing uses the model's own diagnostic prompt to suppress biased generations~\citep{schick2021selfdebiasing}; invariant risk minimization seeks predictors whose optimal decision rule is stable across environments~\citep{arjovsky2019irm}. \auditname does not learn an invariant representation and does not claim that aggregation is unnecessary. Instead, it supplies a paired-control testbed where aggregation-style projections can be audited against the same base-level gap. This is closer in spirit to counterfactual fairness and matched-pair tests, where invariance claims are meaningful only after the intervention and conditioning set are made explicit~\citep{kusner2017counterfactual}.

\paragraph{Positioning.}
The contribution is not a new universal leaderboard, a trained PACT model, or a single-pass replacement for existing guards. It is a protocol-bounded audit-artifact contribution: if matched neutral controls are available, the diagnostic projection can test whether the pressure-specific residual gap remains supported while reporting ordinary F1 only as a non-degradation check. This framing is intentionally narrower than unrestricted robustness leaderboards, but the narrowness is the point: the paper contributes a clean safety-label estimand, base-level inference discipline, and reproducible fail-closed claim gate.
